% =============================================================================
% Ethics Statement (recommended unnumbered for *ACL; not counted in page limit)
% =============================================================================

\ifarxiv
% ---- Full four-paragraph version (arXiv/preprint) ----
\paragraph{Data provenance and licensing.}
SkillCorpus aggregates publicly available \texttt{SKILL.md}
files from GitHub and community marketplaces that link to
GitHub-hosted sources.  Each released skill carries metadata
including its source URL and content hash.  We do not modify
body content or claim authorship.  Every released skill carries
an OSI-approved permissive licence inferred from the source
repository's GitHub \texttt{spdx\_id} and enforced as a Stage-5
ingest filter (Table~\ref{tab:license}, per-reason exclusion
counts in Table~\ref{tab:license-filter}).  Downstream consumers
redistributing the corpus or its derivatives are advised to
perform their own due-diligence audit of \texttt{scripts/}
content and transitive dependencies.

\paragraph{Safety filtering and dual-use.}
The Stage-5 safety hard-gate (\S\ref{sec:classify}) excludes the
most severe categories from the active release set, and the
soft-flag categories are recorded in each skill's metadata for
downstream per-deployment filtering.  We acknowledge the dual-use
risk that a publicly searchable skill corpus could be queried for
adversarial-style content patterns, a residual risk that the
hard-gate and released soft-flag metadata reduce but do not
eliminate.

\paragraph{Benchmark contamination.}
SkillCorpus is community-contributed and pre-dates our benchmark
selection (\S\ref{sec:setup}).  By design the corpus aims for
broad coverage of the skill ecosystem, so released skills overlap
the domains of our evaluation tasks.  These are general procedural
knowledge rather than task-specific solutions, and the per-task
gate injects at most a few topically relevant skills.  We do not
perform task-specific leakage audits beyond the corpus-level
safety and licence filters described above, and consumers using
verifier-backed benchmarks should re-audit for task-specific
contamination when their evaluation regime is sensitive to it.

\else
% ---- Condensed two-paragraph version (AAAI submission, 7-page limit) ----
\paragraph{Data, licensing, and safety.}
SkillCorpus aggregates publicly available \texttt{SKILL.md} files
from GitHub and community marketplaces, recording each skill's
source URL and content hash without modifying its body or claiming
authorship.  Every released skill carries an OSI-approved permissive
licence inferred from the source repository's \texttt{spdx\_id} and
enforced as a Stage-5 gate (Tables~\ref{tab:license}
and~\ref{tab:license-filter}).  The same stage hard-gates the most
severe safety categories and records soft-flags in metadata for
per-deployment filtering.  We acknowledge the residual dual-use risk
that a searchable skill corpus can be queried for adversarial
content, which the gate and flags reduce but do not eliminate.
Downstream consumers redistributing the corpus should audit
\texttt{scripts/} content and transitive dependencies.

\paragraph{Benchmark contamination.}
The corpus is community-contributed and pre-dates our benchmark
selection (\S\ref{sec:setup}).  Its broad coverage overlaps the
domains of our tasks, but the released skills are general procedural
knowledge rather than task-specific solutions, and the per-task gate
injects at most a few topically relevant skills.  We do not perform
task-specific leakage audits beyond the corpus-level filters above,
and consumers using verifier-backed benchmarks should re-audit when
their regime is sensitive to it.
\fi
